% TOP_PROBLEM: {"id": 20001135, "title": "A chain-deficit theorem and exact multidimensional constructions for restricted sums", "category_id": 2, "category": {"id": 2, "name": "combinatorics", "display_name": "Combinatorics", "description": "Counting problems, graph theory, discrete structures.", "slug": "combinatorics", "order_index": 2, "created_at": "2026-07-31T15:26:25.671Z"}, "status": "partially_solved", "problem_number": "AIM-COMBINATORICS-0260", "set_id": 14}
% FIRST_POSTED: 2026-10-01
% ENTRY_KIND: statement_only
% UnsolvedMath ID 20001135; source catalogue code AIM-COMBINATORICS-0260
% !TeX program = lualatex
% Definition and sources only; this file is not an LLM solution attempt.
\documentclass[11pt,a4paper]{article}
\usepackage[margin=25mm]{geometry}
\usepackage{fontspec}
\setmainfont{lmroman10-regular.otf}[BoldFont=lmroman10-bold.otf,
 ItalicFont=lmroman10-italic.otf,BoldItalicFont=lmroman10-bolditalic.otf]
\usepackage{amsmath,amssymb,mathtools}
\usepackage{unicode-math}
\setmathfont{latinmodern-math.otf}
\AtBeginDocument{\let\setminus\smallsetminus}
\usepackage{xurl,hyperref}
\hypersetup{colorlinks=true,urlcolor=blue}
\setcounter{secnumdepth}{0}
\setlength{\parindent}{0pt}
\setlength{\parskip}{5pt}
\setlength{\emergencystretch}{3em}
\begin{document}
\section*{A chain-deficit theorem and exact multidimensional constructions for restricted sums}\label{20001135}
{\small UnsolvedMath ID 20001135 · AIM-COMBINATORICS-0260 · Combinatorics · AIM Workshop Problem Lists}
\subsection{Definitions and mathematical statement}
Let $A\subset\mathbb Z$ have cardinality $n\geq1$, and let $\Gamma\subseteq A\times A$ be an arbitrary set of ordered pairs.
Its restricted sum and difference sets are
\[
 S_\Gamma=\{a+b:(a,b)\in\Gamma\},\qquad
 D_\Gamma=\{a-b:(a,b)\in\Gamma\}.
\]
Pairs with equal coordinates are allowed, and $\Gamma$ need not be invariant under exchanging its coordinates.
Cardinalities count distinct integers, rather than the number of representations or selected pairs.

For each feasible pair of integers $n,m\geq1$, define
\[
 F(n,m)=\max\{|D_\Gamma|:
       A\subset\mathbb Z, |A|=n,
       \Gamma\subseteq A\times A, |S_\Gamma|=m\}.
\]
Determine sharp upper bounds for $F(n,m)$ and describe the integer sets and pair selections attaining, or asymptotically approaching, them.
In particular, identify the structure of $A$ when the number of distinct differences is maximal for a prescribed small number of sums.
The trivial estimates $|D_\Gamma|\leq n(n-1)+1$ and $|D_\Gamma|\leq mn$ only use the number of pairs and the fact that each fixed sum pairs a chosen element with at most one partner.
The question seeks the stronger restrictions coming from simultaneously realizing these pairs within one integer set.
The source's illustrative geometric examples are motivation; they are not imposed as hypotheses or asserted here as exact formulas with unspecified constants.
\subsection{Short English statement}
For a prescribed number of sums on selected pairs of an integer set, maximize the number of differences and describe extremal configurations.
\subsection{Sources}
\begin{itemize}
\item[S1] ULAM.AI and the UnsolvedMath Contributors, supplied problems.json, record 20001135 (AIM-COMBINATORICS-0260), AIM Workshop Problem Lists. Catalogue identity and source transcription; CC BY 4.0.
\url{https://huggingface.co/datasets/ulamai/UnsolvedMath}.
\item[S2] American Institute of Mathematics, Recent trends in additive combinatorics, item 4.4. Original workshop question underlying AIM-COMBINATORICS-0260.
\url{https://aimath.org/WWN/additivecomb/additivecomb.pdf}.
\end{itemize}
\end{document}
